\chapter{Problema, obiettivi e contributi}\label{ch:intro}
Un utente può apprezzare una raccomandazione perché conferma un gusto consolidato oppure perché introduce un artista nuovo. Le due finalità non coincidono: massimizzare la probabilità di ritrovare contenuti già ascoltati può concentrare l'esposizione su una parte ridotta del catalogo; introdurre contenuti molto diversi può aumentare la varietà ma ridurre la pertinenza immediata. Nel dominio musicale, inoltre, il riascolto è un comportamento naturale. L'assenza di novità non costituisce da sola un errore, così come la presenza di un artista mai ascoltato non garantisce un'esperienza positiva.

Il rischio di una \emph{filter bubble} motiva lo studio dell'esplorazione, ma richiede osservazioni sul comportamento e sull'esposizione per essere misurato direttamente. Questo lavoro impiega una definizione più circoscritta di scoperta: una traccia o un artista sono nuovi se non compaiono nella storia disponibile prima della decisione. La novelty così definita è verificabile nei log e distingue l'ampiezza delle proposte dalla capacità di anticipare gli ascolti futuri.

La domanda di ricerca è la seguente: \emph{un modello contestuale con esplorazione può mantenere un'accuratezza di ranking vicina a quella di una baseline personalizzata, aumentando al contempo novelty e diversità delle liste?} Per rispondere occorrono sia un riferimento che sfrutti gli interessi storici, sia strategie che se ne discostino casualmente o in funzione del contesto. La valutazione deve quindi osservare insieme pertinenza e composizione delle proposte, anziché identificare la scoperta con il solo reward. Le policy condividono lo spazio di azioni, ma possono usare segnali diversi.

\section{Scelte progettuali}
Il CMAB consente di rappresentare decisioni ripetute su un catalogo osservabile che cambia, usando storia dell'utente e proprietà dei candidati per orientare l'esplorazione. Il feedback differito è necessario perché il proxy scelto accumula ascolti su più giorni: anticiparlo renderebbe disponibili informazioni non ancora maturate. A ogni inizio di sessione si produce un ranking e si sceglie una sola azione Top-1; la Top-K serve alla valutazione, non rappresenta K azioni premiate.

Il progetto privilegia la tracciabilità della simulazione. Le identità sono condivise fra i moduli; lo stato incrementale esclude gli eventi al timestamp corrente; le statistiche di scala derivano soltanto dai contesti di warm-up. La distinzione fra raccomandazione track-level e rilevanza artist-level è esplicita: le Top-K tracce vengono convertite in una sequenza di artisti distinti per Precision, Recall e NDCG. Questa scelta evita che ripetere lo stesso artista generi più rilevanze indipendenti.

\section{Contributo e perimetro delle conclusioni}
I contributi sono un simulatore causale con feedback differito, un confronto riproducibile fra cinque famiglie di policy, una selezione congelata prima del test e un'analisi dell'incertezza che mantiene raggruppate le decisioni dello stesso utente. I risultati includono anche esempi musicali autentici, utili per interpretare le liste al di là delle medie.

In un sistema online l'esplorazione potrebbe cambiare le esposizioni e raccogliere nuove risposte. Qui la sequenza Last.fm rimane fissa: raccomandare un artista non ne modifica gli ascolti futuri. Si studia quindi la traiettoria decisionale del bandit rispetto a un proxy osservato, non la sua capacità causale di cambiare i gusti. Questa distinzione motiva la lettura congiunta delle metriche del Capitolo~\ref{ch:metodo} e delimita le conclusioni.
